\scanpage{29}
\textcircled{2} Cas où $M$ est un $\mathbb Z_p$-module

Notons d'abord qu'il existe des $k$ arbitrairem\textsuperscript{t} grands tels que $c_{i,k-i}\in\mathbb Z_p^{\unsure{\times|*}}$ i.e. $c_{i,k-i}\not\equiv0\ (p)$. Posant $k-i=j$, il suffit de prendre $j=p^\varepsilon$, avec $\varepsilon$ tel que $p^\varepsilon>i$, donc $k=i+p^\varepsilon$. Cela montre l'unicité de $\xi$, car si $k_0$ est tel que $c_{i,k_0-i}\in\mathbb Z_p^\times$, on aura
\[
\deleted{\unsure{(2)}}\qquad \xi=c_{i,k_0-i}^{-1}\xi_{k_0}
\]
Considérons alors le système des
\[
\deleted{\unsure{(3)}}\qquad \xi'_k=c_{i,k-i}\xi.\quad,
\]
on voit de suite que \deleted{\unsure{ce}} $(\xi'_k)_{k\geq N}$ satisfait \deleted{\unsure{aux hyp|analogue}} les conditions (1)\unsure{,|{}} (où on substitue $\xi'_k$ à $\xi_k$), et $\xi'_{k_0}=\xi_{k_0}$, donc posons
\[
\unsure{\eta|y}_k=\xi'_k-\xi_k,
\]
\deleted{les} $(\eta_k)_{k\geq N}$ satisfait encore (1), et de plus $\eta_{k_0}=0$. Il suffit de prouver alors $\eta_k=0$, \deleted{donc} $\forall k\geq N$. Donc on est ramené au

\uline{Corollaire} Si $(\xi_k)_{k\geq N}$ \deleted{est} satisfait (1), et \inserted{si} de plus \deleted{\unsure{$(1\frac12)$}} on \unsure{donné|donne} un $k_0\geq N$ tel que
\[
\xi_{k_0}=0\tag*{$(1')$}
\]
\[
c_{i,k_0-i}\not\equiv0\ (p)\tag*{(3)}
\]
alors $\qquad \xi_k=0\qquad\forall k\geq N$.
% archive sheet number: 28

